\documentclass[11pt,a4paper]{article}

% ================ Encodage et langue ================
\usepackage[utf8]{inputenc}
\usepackage[T1]{fontenc}
\usepackage[french]{babel}
\usepackage{lmodern}
\usepackage{textcomp}

% ================ Mise en page ================
\usepackage[margin=2.2cm]{geometry}
\usepackage{parskip}
\setlength{\parindent}{0pt}
\setlength{\parskip}{0.6em}

% ================ Couleurs et liens ================
\usepackage[table]{xcolor}
\definecolor{primary}{HTML}{0B5394}
\definecolor{accent}{HTML}{1F77B4}
\definecolor{light}{HTML}{F2F2F2}
\definecolor{codebg}{HTML}{F6F8FA}
\definecolor{rulec}{HTML}{D0D7DE}

\usepackage[colorlinks=true, linkcolor=primary, urlcolor=primary, citecolor=primary]{hyperref}

% ================ Tableaux ================
\usepackage{longtable, booktabs, array, tabularx}
\renewcommand{\arraystretch}{1.25}
\arrayrulecolor{rulec}

% ================ Listes ================
\usepackage{enumitem}
\setlist{itemsep=2pt, topsep=4pt}

% ================ Code et citations ================
\usepackage{fancyvrb}
\usepackage{framed}
\definecolor{shadecolor}{named}{light}

% Encadrés pour les blocs > (quotes)
\usepackage{mdframed}
\newmdenv[
  linecolor=primary,
  linewidth=2pt,
  topline=false, bottomline=false, rightline=false,
  innerleftmargin=10pt, innerrightmargin=8pt,
  innertopmargin=6pt, innerbottommargin=6pt,
  backgroundcolor=codebg,
  skipabove=8pt, skipbelow=8pt
]{quotebox}
\renewenvironment{quote}{\begin{quotebox}}{\end{quotebox}}

% ================ Titres ================
\usepackage{titlesec}
\titleformat{\section}{\Large\bfseries\color{primary}}{\thesection.}{0.6em}{}
\titleformat{\subsection}{\large\bfseries\color{accent}}{\thesubsection}{0.6em}{}
\titleformat{\subsubsection}{\normalsize\bfseries}{\thesubsubsection}{0.6em}{}

% ================ En-têtes et pieds de page ================
\usepackage{fancyhdr}
\pagestyle{fancy}
\fancyhf{}
\fancyhead[L]{\small\color{gray}ESPRIT School of Business --- M1 BA --- Projet Intégré}
\fancyhead[R]{\small\color{gray}\nouppercase{\leftmark}}
\fancyfoot[C]{\small\thepage}
\renewcommand{\headrulewidth}{0.3pt}
\renewcommand{\footrulewidth}{0pt}

% ================ Pandoc helpers ================
\providecommand{\tightlist}{\setlength{\itemsep}{0pt}\setlength{\parskip}{0pt}}
\providecommand{\passthrough}[1]{#1}

\usepackage{calc}
\providecommand{\real}[1]{#1}


% ================ Métadonnées ================
\title{\color{primary}\textbf{Timeline du Projet --- 4 Semaines}}
\author{Aymen Ben Brik}
\date{2026}

\begin{document}

\begin{titlepage}
\centering
\vspace*{1.5cm}

{\Large\bfseries ESPRIT School of Business}\\[0.3cm]
{\large Master 1 --- Business Analytics (M1 BA)}\\[2cm]

{\color{gray}\large Projet Intégré}\\[0.3cm]
{\color{gray}\large Analyse et Prédiction du Churn Client}\\[2.2cm]

\rule{0.85\linewidth}{1.2pt}\\[0.6cm]
{\Huge\bfseries\color{primary} Timeline du Projet --- 4 Semaines\par}
\rule{0.85\linewidth}{1.2pt}\\[2cm]

{\large Tuteur du projet}\\[0.2cm]
{\Large\bfseries Aymen Ben Brik}\\[1cm]
{\large 2026}

\vfill
{\color{gray}\small Document à destination des étudiants}
\end{titlepage}

\tableofcontents
\newpage

\section{Timeline du Projet --- 4
Semaines}\label{timeline-du-projet-4-semaines}

La durée totale du projet est de \textbf{4 semaines (28 jours)}. Le
planning ci-dessous est indicatif : adaptez-le à la composition et à la
disponibilité de votre équipe, mais respectez les jalons hebdomadaires
qui correspondent aux points de suivi avec l'encadrant.

\subsection{Vue d'ensemble}\label{vue-densemble}

{\def\LTcaptype{none} % do not increment counter
\begin{longtable}[]{@{}
  >{\raggedright\arraybackslash}p{(\linewidth - 6\tabcolsep) * \real{0.2500}}
  >{\raggedright\arraybackslash}p{(\linewidth - 6\tabcolsep) * \real{0.2500}}
  >{\raggedright\arraybackslash}p{(\linewidth - 6\tabcolsep) * \real{0.2500}}
  >{\raggedright\arraybackslash}p{(\linewidth - 6\tabcolsep) * \real{0.2500}}@{}}
\toprule\noalign{}
\begin{minipage}[b]{\linewidth}\raggedright
Semaine
\end{minipage} & \begin{minipage}[b]{\linewidth}\raggedright
Phase
\end{minipage} & \begin{minipage}[b]{\linewidth}\raggedright
Objectif principal
\end{minipage} & \begin{minipage}[b]{\linewidth}\raggedright
Livrable de fin de semaine
\end{minipage} \\
\midrule\noalign{}
\endhead
\bottomrule\noalign{}
\endlastfoot
\textbf{S1} & Cadrage \& Exploration & Comprendre les données, mettre en
place l'environnement et l'ETL & Pipeline ETL fonctionnel + rapport
d'exploration \\
\textbf{S2} & BI \& Power BI & Modélisation dimensionnelle et tableaux
de bord & Data Warehouse + rapports Power BI \\
\textbf{S3} & Machine Learning & Entraînement et évaluation des modèles
de churn & Modèle sauvegardé + notebook d'analyse \\
\textbf{S4} & Web \& Livraison & Application web, déploiement, rapport,
présentation & Démo en ligne + rapport PDF + présentation \\
\end{longtable}
}

\begin{center}\rule{0.5\linewidth}{0.5pt}\end{center}

\subsection{Semaine 1 --- Cadrage et Exploration (J1 →
J7)}\label{semaine-1-cadrage-et-exploration-j1-j7}

\subsubsection{Objectifs}\label{objectifs}

\begin{itemize}
\tightlist
\item
  Constituer l'équipe et définir les rôles.
\item
  Comprendre le contexte métier et la définition du churn.
\item
  Profiler les données et identifier les problèmes de qualité.
\item
  Concevoir et implémenter le pipeline ETL.
\end{itemize}

\subsubsection{Jour par jour}\label{jour-par-jour}

{\def\LTcaptype{none} % do not increment counter
\begin{longtable}[]{@{}
  >{\raggedright\arraybackslash}p{(\linewidth - 2\tabcolsep) * \real{0.5000}}
  >{\raggedright\arraybackslash}p{(\linewidth - 2\tabcolsep) * \real{0.5000}}@{}}
\toprule\noalign{}
\begin{minipage}[b]{\linewidth}\raggedright
Jour
\end{minipage} & \begin{minipage}[b]{\linewidth}\raggedright
Activité
\end{minipage} \\
\midrule\noalign{}
\endhead
\bottomrule\noalign{}
\endlastfoot
J1 & Lecture des documents du projet, répartition des rôles, création du
dépôt GitHub, mise en place de l'environnement (Python, venv,
dépendances). \\
J2 & Exploration des fichiers : volumétrie, types, valeurs manquantes,
distributions. Premier notebook d'EDA (Exploratory Data Analysis). \\
J3 & Analyse des tables de dimensions, étude des codes et libellés,
recensement des incohérences. \\
J4 & Définition formelle de la variable cible (churn), justification du
choix. Liste des features candidates. \\
J5 & Conception du pipeline ETL : sources, transformations, cibles.
Diagramme. \\
J6 & Implémentation du pipeline ETL (extraction + nettoyage +
jointures). \\
J7 & Chargement vers la base de données analytique. Tests de cohérence.
\textbf{Jalon S1.} \\
\end{longtable}
}

\subsubsection{Livrable de la semaine}\label{livrable-de-la-semaine}

\begin{itemize}
\tightlist
\item
  Dépôt GitHub initialisé avec README, structure de dossiers,
  \texttt{.gitignore}.
\item
  Notebook d'exploration documenté.
\item
  Pipeline ETL exécutable et reproductible.
\item
  Schéma de l'architecture cible (data flow).
\end{itemize}

\begin{center}\rule{0.5\linewidth}{0.5pt}\end{center}

\subsection{Semaine 2 --- BI et Power BI (J8 →
J14)}\label{semaine-2-bi-et-power-bi-j8-j14}

\subsubsection{Objectifs}\label{objectifs-1}

\begin{itemize}
\tightlist
\item
  Concevoir un modèle dimensionnel en étoile.
\item
  Définir et calculer les KPIs métier.
\item
  Produire les rapports Power BI.
\end{itemize}

\subsubsection{Jour par jour}\label{jour-par-jour-1}

{\def\LTcaptype{none} % do not increment counter
\begin{longtable}[]{@{}
  >{\raggedright\arraybackslash}p{(\linewidth - 2\tabcolsep) * \real{0.5000}}
  >{\raggedright\arraybackslash}p{(\linewidth - 2\tabcolsep) * \real{0.5000}}@{}}
\toprule\noalign{}
\begin{minipage}[b]{\linewidth}\raggedright
Jour
\end{minipage} & \begin{minipage}[b]{\linewidth}\raggedright
Activité
\end{minipage} \\
\midrule\noalign{}
\endhead
\bottomrule\noalign{}
\endlastfoot
J8 & Conception du modèle dimensionnel : table de faits, dimensions,
granularité. \\
J9 & Création des tables dimensionnelles et de la table de faits dans le
Data Warehouse. \\
J10 & Définition des KPIs : taux de churn global, par segment, par
ancienneté, par produit, etc. \\
J11 & Connexion de Power BI à la source de données. Préparation des
mesures DAX. \\
J12 & Construction du rapport Power BI --- page 1 (vue d'ensemble). \\
J13 & Construction du rapport Power BI --- pages d'analyse
(segmentation, churn, produits). \\
J14 & Mise au point graphique, interactivité, publication. \textbf{Jalon
S2.} \\
\end{longtable}
}

\subsubsection{Livrable de la semaine}\label{livrable-de-la-semaine-1}

\begin{itemize}
\tightlist
\item
  Schéma dimensionnel documenté.
\item
  Data Warehouse peuplé.
\item
  Fichier \texttt{.pbix} (ou lien Power BI Service) avec au minimum 3
  pages d'analyse.
\item
  Liste des KPIs documentée dans le dépôt.
\end{itemize}

\begin{center}\rule{0.5\linewidth}{0.5pt}\end{center}

\subsection{Semaine 3 --- Machine Learning (J15 →
J21)}\label{semaine-3-machine-learning-j15-j21}

\subsubsection{Objectifs}\label{objectifs-2}

\begin{itemize}
\tightlist
\item
  Préparer le jeu de données ML.
\item
  Entraîner et comparer plusieurs modèles.
\item
  Choisir et sauvegarder le modèle final.
\end{itemize}

\subsubsection{Jour par jour}\label{jour-par-jour-2}

{\def\LTcaptype{none} % do not increment counter
\begin{longtable}[]{@{}
  >{\raggedright\arraybackslash}p{(\linewidth - 2\tabcolsep) * \real{0.5000}}
  >{\raggedright\arraybackslash}p{(\linewidth - 2\tabcolsep) * \real{0.5000}}@{}}
\toprule\noalign{}
\begin{minipage}[b]{\linewidth}\raggedright
Jour
\end{minipage} & \begin{minipage}[b]{\linewidth}\raggedright
Activité
\end{minipage} \\
\midrule\noalign{}
\endhead
\bottomrule\noalign{}
\endlastfoot
J15 & Préparation du dataset ML depuis le DWH : sélection des features,
encodage, traitement du déséquilibre. \\
J16 & Split train/test, baseline (régression logistique). Premières
métriques. \\
J17 & Entraînement de modèles avancés (Random Forest, XGBoost,
LightGBM). Comparaison. \\
J18 & Optimisation des hyperparamètres (Grid Search ou Optuna) pour le
meilleur candidat. \\
J19 & Évaluation finale : accuracy, precision, recall, F1, ROC-AUC,
matrice de confusion. \\
J20 & Interprétabilité : importance des features, SHAP values. \\
J21 & Sauvegarde du modèle (\texttt{.pkl} ou \texttt{.joblib}).
Documentation des choix. \textbf{Jalon S3.} \\
\end{longtable}
}

\subsubsection{Livrable de la semaine}\label{livrable-de-la-semaine-2}

\begin{itemize}
\tightlist
\item
  Notebook ML complet (préparation → entraînement → évaluation →
  interprétation).
\item
  Modèle sauvegardé.
\item
  Tableau comparatif des modèles testés.
\end{itemize}

\begin{center}\rule{0.5\linewidth}{0.5pt}\end{center}

\subsection{Semaine 4 --- Interface Web, Déploiement, Livraison (J22 →
J28)}\label{semaine-4-interface-web-duxe9ploiement-livraison-j22-j28}

\subsubsection{Objectifs}\label{objectifs-3}

\begin{itemize}
\tightlist
\item
  Construire et déployer l'application web.
\item
  Finaliser le rapport et la présentation.
\end{itemize}

\subsubsection{Jour par jour}\label{jour-par-jour-3}

{\def\LTcaptype{none} % do not increment counter
\begin{longtable}[]{@{}
  >{\raggedright\arraybackslash}p{(\linewidth - 2\tabcolsep) * \real{0.5000}}
  >{\raggedright\arraybackslash}p{(\linewidth - 2\tabcolsep) * \real{0.5000}}@{}}
\toprule\noalign{}
\begin{minipage}[b]{\linewidth}\raggedright
Jour
\end{minipage} & \begin{minipage}[b]{\linewidth}\raggedright
Activité
\end{minipage} \\
\midrule\noalign{}
\endhead
\bottomrule\noalign{}
\endlastfoot
J22 & Maquette de l'application web. Choix de la stack (Streamlit /
Flask / FastAPI + frontend). \\
J23 & Développement --- chargement du modèle, formulaire de prédiction
client. \\
J24 & Développement --- page ``comptes à risque'' + KPIs de synthèse. \\
J25 & Déploiement sur une plateforme cloud (Streamlit Cloud, Render,
Hugging Face Spaces, Railway, etc.). Tests. \\
J26 & Rédaction du rapport PDF (toutes les parties). \\
J27 & Préparation de la présentation (slides + script +
démonstration). \\
J28 & Répétition de la présentation. Finalisation du dépôt GitHub.
\textbf{Soutenance finale.} \\
\end{longtable}
}

\subsubsection{Livrable final}\label{livrable-final}

\begin{itemize}
\tightlist
\item
  Application web accessible en ligne via une URL publique.
\item
  Dépôt GitHub finalisé avec README détaillé (instructions, captures
  d'écran, lien démo).
\item
  Rapport PDF (25--40 pages).
\item
  Slides de présentation.
\end{itemize}

\begin{center}\rule{0.5\linewidth}{0.5pt}\end{center}

\subsection{Recommandations
transverses}\label{recommandations-transverses}

\begin{itemize}
\tightlist
\item
  \textbf{Commits réguliers} : chaque membre commit au moins
  quotidiennement sur sa branche.
\item
  \textbf{Branches Git} : utilisez \texttt{main} pour le code stable et
  des branches \texttt{feature/...} par tâche.
\item
  \textbf{Documentation au fil de l'eau} : ne laissez pas la rédaction
  du rapport pour la dernière semaine. Rédigez les sections au fur et à
  mesure que vous terminez les phases.
\item
  \textbf{Tests} : prévoyez du temps de buffer en fin de semaine 4 ---
  déployer pour la première fois prend souvent plus de temps que prévu.
\end{itemize}

\end{document}
